\noindent\textbf{Prompt}\par\smallskip

\begin{tcolorbox}[
  breakable, enhanced,
  title={System},
  colback=SystemBg, colframe=SystemFrame,
  coltitle=white, fonttitle=\bfseries\small,
  arc=2mm, boxrule=1pt,
  before upper={\small\setlength{\parindent}{0pt}\setlength{\parskip}{2pt}},
]
Your input fields are:\\{}
1. `system\_description` (str): System description\\{}
2. `simulator\_code` (str): Simulator code\\{}
3. `performance\_metrics` (str): Performance metrics\\{}
Your output fields are:\\{}
1. `feedback` (ProvideFeedbackOutput): Improvement feedback\\{}
All interactions will be structured in the following way, with the appropriate values filled in.
\smallskip
\par\noindent [[ \#\# system\_description \#\# ]]\\
\noindent
\{system\_description\}
\smallskip
\par\noindent [[ \#\# simulator\_code \#\# ]]\\
\noindent
\{simulator\_code\}
\smallskip
\par\noindent [[ \#\# performance\_metrics \#\# ]]\\
\noindent
\{performance\_metrics\}
\smallskip
\par\noindent [[ \#\# feedback \#\# ]]\\
\noindent
\{feedback\}        \# note: the value you produce must adhere to the JSON schema: \{"type": "object", "\$defs": \{"Issue": \{"type": "object", "properties": \{"description": \{"type": "string", "description": "Issue description", "title": "Description"\}, "location": \{"anyOf": [\{"type": "string"\}, \{"type": "null"\}], "default": null, "description": "Code location", "title": "Location"\}, "severity": \{"type": "string", "description": "Severity level", "enum": ["critical", "major", "minor", "suggestion"], "title": "Severity"\}, "suggestion": \{"type": "string", "description": "Fix suggestion", "title": "Suggestion"\}\}, "required": ["description", "severity", "suggestion"], "title": "Issue"\}\}, "properties": \{"issues": \{"type": "array", "description": "Issues with fixes", "items": \{"\$ref": "\#/\$defs/Issue"\}, "maxItems": 2, "title": "Issues"\}, "main\_diagnosis": \{"type": "string", "description": "Main issue and location", "title": "Main Diagnosis"\}\}, "required": ["main\_diagnosis", "issues"], "title": "ProvideFeedbackOutput"\}
\smallskip
\par\noindent [[ \#\# completed \#\# ]]\\
\noindent
In adhering to this structure, your objective is: \\{}
        Analyze evaluation metrics to diagnose failure modes, then provide code improvement feedback.
\smallskip
\end{tcolorbox}
\vspace{0.35em}

\begin{tcolorbox}[
  breakable, enhanced,
  title={User},
  colback=UserBg, colframe=UserFrame,
  coltitle=white, fonttitle=\bfseries\small,
  arc=2mm, boxrule=1pt,
  before upper={\small\setlength{\parindent}{0pt}\setlength{\parskip}{2pt}},
]
\par\noindent [[ \#\# system\_description \#\# ]]\\
\noindent
\# HODGKIN-HUXLEY NEURON SIMULATOR WITH MISSING ION CHANNELS
\par\smallskip
\#\# OBJECTIVE\\{}
Extend an existing Hodgkin-Huxley neuron simulator by identifying and implementing\\{}
missing ion channels necessary to reduce discrepancies between simulations and\\{}
experimental voltage recordings. The goal is to improve agreement across multiple\\{}
electrophysiological summary statistics while maintaining model parsimony.
\par\smallskip
\#\# BASE MODEL\\{}
Classic Hodgkin-Huxley formulation with standard three-gating-variable structure:\\{}
- **Leak current**: conductance g\_leak, reversal potential E\_leak\\{}
- **Sodium (Na$^{+}$) current**: conductance gNa, gating variables m (activation) and h (inactivation), reversal E\_Na\\{}
- **Potassium (K$^{+}$) current**: conductance gK, gating variable n (activation), reversal E\_K\\{}
- **Dynamics**: First-order gating kinetics with voltage-dependent time constants and steady states\\{}
- **Input**: Externally applied current I\_inj(t) with optional stochastic components\\{}
- **Stochasticity**: Optional random seed for reproducible simulations
\par\smallskip
\#\# EXTENSIBILITY\\{}
Two pre-allocated channel slots (X1, X2) are available for additional mechanisms.\\{}
Each slot provides two tunable parameters (see signature\_description for details).\\{}
Channels beyond these two slots must use fixed parameters, not inferred ones.
\par\smallskip
\#\# EVALUATION METRICS\\{}
Model performance assessed on held-out voltage recordings via these summary statistics:\\{}
1. Number of spikes during stimulation\\{}
2. Mean resting potential (pre-stimulation)\\{}
3. Standard deviation of resting potential\\{}
4. Mean voltage during stimulation\\{}
5. Variance of voltage\\{}
6. Skewness of voltage\\{}
7. Kurtosis of voltage
\par\smallskip
\#\# DATA CONTEXT\\{}
Intracellular voltage recordings from a single neuron exhibiting tonic spiking\\{}
behavior across different stimulus intensities.
\par\smallskip
**Important data characteristics**:\\{}
- No burst firing patterns (no clusters of rapid spikes separated by quiescence)\\{}
- No prolonged or sustained high-frequency firing (no sustained high-frequency firing)\\{}
- Regular, tonic spiking activity during periods of input current stimulation\\{}
- Non-spiking (quiescent) behavior in the absence of input current\\{}
- Evenly-spaced action potentials without burst patterns during the activation period
\par\smallskip
Therefore, avoid adding channels specifically designed to produce bursting or sustained\\{}
high-frequency dynamics.
\smallskip
\par\noindent [[ \#\# simulator\_code \#\# ]]\\
\noindent
\begin{lstlisting}[language=Python,style=pythonstyle]
import torch
import torch.nn as nn


class DiscoveredSimulator(nn.Module):
 def __init__(self):
  super(DiscoveredSimulator, self).__init__()
  return

 def forward(
  self,
  init_voltage: float,
  input_current: torch.Tensor,
  dt: float,
  t: torch.Tensor,
  params: torch.Tensor,
  seed=None,
 ):
  """
  Hodgkin-Huxley neuron extended with:
   1. M-current (I_KM): slow non-inactivating K+ current using X1 slot
    - Provides spike-frequency adaptation, stabilises tonic firing
    - Uses voltage-independent tau (more stable during spikes)
    - gbar_KM = params[:,6], V_half_p = -params[:,8], tau_p = params[:,9]

   2. Ih current (HCN channel): hyperpolarisation-activated cation current using X2 slot
    - Opens at hyperpolarised voltages (~-80 mV and below)
    - Reversal E_h ~= -30 mV (mixed Na+/K+), depolarising at rest
    - Corrects resting potential mean, SD, and voltage distribution shape
    - Fixed kinetics (half-activation -80 mV, tau 200 ms); only gbar_Ih inferred
    - gbar_Ih = params[:,7], all other Ih parameters are fixed

  Args:
   init_voltage: torch.Tensor: (batch_size,)            initial voltage [mV]
   input_current: torch.Tensor: (batch_size, time_steps) injected current [uA/cm2]
   dt: float                                              time step [ms]
   t: torch.Tensor: (time_steps,)                        time array [ms]
   params: torch.Tensor: (batch_size, 10)                biophysical parameters
   seed: optional random seed

  Returns:
   V: torch.Tensor: (batch_size, time_steps)             voltage traces [mV]
  """
  device = params.device

  # Set up random generator
  if seed is not None:
   generator = torch.Generator(device=device)
   generator.manual_seed(seed)
  else:
   generator = torch.Generator(device=device)

  batch_size = params.shape[0]   # int
  time_steps = t.shape[0]        # int

  # -- Base HH parameters ----------------------------------------------------
  gbar_Na   = params[:, 0].float()   # (batch_size,)  mS/cm^2
  gbar_K    = params[:, 1].float()   # (batch_size,)  mS/cm^2
  g_leak    = params[:, 2].float()   # (batch_size,)  mS/cm^2
  E_leak    = -params[:, 3].float()  # (batch_size,)  mV  (negative applied)
  Vt        = -params[:, 4].float()  # (batch_size,)  mV  (negative applied)
  nois_fact = params[:, 5].float()   # (batch_size,)  unitless

  # -- X1 slot: M-current (I_KM, slow non-inactivating K+) -----------------
  # Physiological rationale:
  #   M-current activates slowly at subthreshold depolarised voltages (~-35 mV),
  #   providing an outward K+ current that limits repetitive firing rate and
  #   produces spike-frequency adaptation without bursting or quiescence.
  # Parameter mapping:
  #   gbar_KM   = params[:,6] in [1e-4, 10]  mS/cm^2
  #   V_half_p  = -params[:,8] in [-150, -1e-4] mV  (typically ~-35 mV)
  #   tau_p     =  params[:,9] in [1e-4, 3000]  ms  (typically 50-500 ms)
  # Note: voltage-independent tau avoids instability during fast action potentials
  gbar_KM  = params[:, 6].float()   # (batch_size,)  mS/cm^2
  V_half_p = -params[:, 8].float()  # (batch_size,)  mV
  tau_p    =  params[:, 9].float()  # (batch_size,)  ms --- voltage-independent for stability

  # -- X2 slot: Ih current (HCN, hyperpolarisation-activated cation) --------
  # Physiological rationale:
  #   Ih activates at hyperpolarised voltages (< -60 mV), carries inward
  #   Na+/K+ current with reversal ~-30 mV, depolarising the cell toward
  #   rest. This sets a tonic inward "sag" current that:
  #     - Raises mean resting potential slightly above pure K+ equilibrium
  #     - Increases resting SD (voltage fluctuations due to channel noise)
  #     - Shapes subthreshold distribution (skewness, kurtosis)
  #   It does NOT promote bursting --- it acts as a stabilising "pacemaker"
  #   current opposing excessive hyperpolarisation.
  # Parameter mapping:
  #   gbar_Ih = params[:,7] in [1e-4, 120] mS/cm^2 --- only inferred parameter
  #   All kinetic parameters fixed from literature
  gbar_Ih = params[:, 7].float()   # (batch_size,)  mS/cm^2

  # Fixed Ih kinetic constants (from Magee 1998, Koch 1999)
  E_h          = -30.0   # mV  --- mixed Na+/K+ cation reversal
  V_half_r     = -80.0   # mV  --- half-activation (hyperpolarisation-activated)
  k_r          =  7.0    # mV  --- activation slope (negative: opens on hyperpol.)
  tau_r_fixed  = 200.0   # ms  --- slow, voltage-independent time constant

  tstep = float(dt)

  # -- Fixed biophysical constants -------------------------------------------
  nois_fact_obs = 0.0
  C    = 1.0    # uF/cm^2
  E_Na = 53.0   # mV
  E_K  = -107.0 # mV

  # -- Numerical helpers -----------------------------------------------------
  def Exp(z):
   # Numerically safe exponential --- prevent overflow at very negative z
   return torch.where(
    z < -5e2,
    torch.exp(torch.full_like(z, -5e2)),
    torch.exp(z)
   )  # same shape as z

  def efun(z):
   # Handles 0/0 limit of z / (exp(z) - 1) near z=0
   return torch.where(
    torch.abs(z) < 1e-4,
    1 - z / 2,
    z / (Exp(z) - 1)
   )  # same shape as z

  # -- Standard HH channel kinetics -----------------------------------------
  def alpha_m(x):  # (batch_size,) -> (batch_size,)
   v1 = x - Vt - 13.0
   return 0.32 * efun(-0.25 * v1) / 0.25

  def beta_m(x):   # (batch_size,) -> (batch_size,)
   v1 = x - Vt - 40.0
   return 0.28 * efun(0.2 * v1) / 0.2

  def alpha_h(x):  # (batch_size,) -> (batch_size,)
   v1 = x - Vt - 17.0
   return 0.128 * Exp(-v1 / 18.0)

  def beta_h(x):   # (batch_size,) -> (batch_size,)
   v1 = x - Vt - 40.0
   return 4.0 / (1 + Exp(-0.2 * v1))

  def alpha_n(x):  # (batch_size,) -> (batch_size,)
   v1 = x - Vt - 15.0
   return 0.032 * efun(-0.2 * v1) / 0.2

  def beta_n(x):   # (batch_size,) -> (batch_size,)
   v1 = x - Vt - 10.0
   return 0.5 * Exp(-v1 / 40)

  def tau_x(alpha, beta):  # (batch_size,), (batch_size,) -> (batch_size,)
   return 1.0 / (alpha + beta)

  def inf_x(alpha, beta):  # (batch_size,), (batch_size,) -> (batch_size,)
   return alpha / (alpha + beta)

  # -- M-current (I_KM) gating: Boltzmann steady-state, constant tau --------
  # p_inf = 1 / (1 + exp(-(V - V_half_p) / 10))
  # Slope k=10 mV is standard for cortical M-current (Wang 1998)
  def inf_p(x):  # (batch_size,) -> (batch_size,)
   return 1.0 / (1.0 + Exp(-(x - V_half_p) / 10.0))

  # Voltage-independent tau_p avoids spurious fast tracking during spikes.
  # tau_p is directly the inferred parameter (batch_size,), already positive.

  # -- Ih gating: hyperpolarisation-activated Boltzmann, fixed tau ----------
  # r_inf = 1 / (1 + exp((V - V_half_r) / k_r))
  # Note positive sign in exponent: more open at V << V_half_r = -80 mV
  def inf_r(x):  # (batch_size,) -> (batch_size,)
   return 1.0 / (1.0 + Exp((x - V_half_r) / k_r))

  # tau_r is fixed scalar (200 ms) --- no per-batch variation needed

  # -- State variable arrays -------------------------------------------------
  V = torch.zeros((batch_size, time_steps), device=device)  # (batch_size, time_steps)
  n = torch.zeros((batch_size, time_steps), device=device)  # (batch_size, time_steps)
  m = torch.zeros((batch_size, time_steps), device=device)  # (batch_size, time_steps)
  h = torch.zeros((batch_size, time_steps), device=device)  # (batch_size, time_steps)
  p = torch.zeros((batch_size, time_steps), device=device)  # (batch_size, time_steps) M-current gate
  r = torch.zeros((batch_size, time_steps), device=device)  # (batch_size, time_steps) Ih gate

  # -- Initial conditions (steady state at initial voltage) ------------------
  V_init = init_voltage.to(device)  # (batch_size,)
  V[:, 0] = V_init                                               # (batch_size,)
  n[:, 0] = inf_x(alpha_n(V[:, 0]), beta_n(V[:, 0]))            # (batch_size,)
  m[:, 0] = inf_x(alpha_m(V[:, 0]), beta_m(V[:, 0]))            # (batch_size,)
  h[:, 0] = inf_x(alpha_h(V[:, 0]), beta_h(V[:, 0]))            # (batch_size,)
  p[:, 0] = inf_p(V[:, 0])                                       # (batch_size,)
  r[:, 0] = inf_r(V[:, 0])                                       # (batch_size,)

  # -- Exponential Euler time-integration loop -------------------------------
  for i in range(1, time_steps):
   V_prev = V[:, i - 1]  # (batch_size,)

   # Standard HH gating rates
   a_m, b_m = alpha_m(V_prev), beta_m(V_prev)   # (batch_size,), (batch_size,)
   a_h, b_h = alpha_h(V_prev), beta_h(V_prev)   # (batch_size,), (batch_size,)
   a_n, b_n = alpha_n(V_prev), beta_n(V_prev)   # (batch_size,), (batch_size,)

   # M-current gate steady-state (tau is voltage-independent -> use tau_p directly)
   inf_p_val = inf_p(V_prev)    # (batch_size,)

   # Ih gate steady-state (tau is fixed scalar)
   inf_r_val = inf_r(V_prev)    # (batch_size,)

   # Effective conductances at previous time step
   g_Na_eff  = (m[:, i-1] ** 3) * gbar_Na * h[:, i-1]  # (batch_size,)
   g_K_eff   = (n[:, i-1] ** 4) * gbar_K                # (batch_size,)
   g_KM_eff  = gbar_KM * p[:, i-1]                      # (batch_size,)
   g_Ih_eff  = gbar_Ih * r[:, i-1]                      # (batch_size,)

   # Effective inverse membrane time constant (total conductance / C)
   tau_V_inv = (
    g_Na_eff
    + g_K_eff
    + g_leak
    + g_KM_eff   # M-current: outward K+, raises conductance
    + g_Ih_eff   # Ih: inward mixed cation, raises conductance
   ) / C  # (batch_size,)

   # Voltage steady-state numerator (weighted reversal potentials + input)
   V_inf = (
    g_Na_eff * E_Na
    + g_K_eff  * E_K
    + g_leak   * E_leak
    + g_KM_eff * E_K    # M-current drives V toward E_K (hyperpolarising)
    + g_Ih_eff * E_h    # Ih drives V toward E_h ~ -30 mV (depolarising at rest)
    + input_current[:, i-1]
    + nois_fact * torch.randn(batch_size, generator=generator, device=device) / (tstep ** 0.5)
   ) / (tau_V_inv * C)  # (batch_size,)

   # Exponential Euler update for voltage
   V[:, i] = V_inf + (V_prev - V_inf) * Exp(-tstep * tau_V_inv)          # (batch_size,)

   # Standard HH gating variable updates
   n[:, i] = inf_x(a_n, b_n) + (n[:, i-1] - inf_x(a_n, b_n)) * Exp(-tstep / tau_x(a_n, b_n))  # (batch_size,)
   m[:, i] = inf_x(a_m, b_m) + (m[:, i-1] - inf_x(a_m, b_m)) * Exp(-tstep / tau_x(a_m, b_m))  # (batch_size,)
   h[:, i] = inf_x(a_h, b_h) + (h[:, i-1] - inf_x(a_h, b_h)) * Exp(-tstep / tau_x(a_h, b_h))  # (batch_size,)

   # M-current gate update (voltage-independent tau -> stable during spikes)
   p[:, i] = inf_p_val + (p[:, i-1] - inf_p_val) * Exp(-tstep / tau_p)  # (batch_size,)

   # Ih gate update (fixed scalar tau)
   r[:, i] = inf_r_val + (r[:, i-1] - inf_r_val) * Exp(
    torch.full_like(V_prev, -tstep / tau_r_fixed)
   )  # (batch_size,)

  # Return voltage traces (+ optional observation noise, currently zero)
  return V + nois_fact_obs * torch.randn(
   batch_size, time_steps, generator=generator, device=device
  )  # (batch_size, time_steps)
\end{lstlisting}
\smallskip
\par\noindent [[ \#\# performance\_metrics \#\# ]]\\
\noindent
Metrics: neg\_avg\_log\_marginal\_NLE: 35.3, neg\_log\_marginal\_NLE: 353, timeout\_evidence\_estimation: 0
\par\smallskip
Respond with the corresponding output fields, starting with the field `[[ \#\# feedback \#\# ]]` (must be formatted as a valid Python ProvideFeedbackOutput), and then ending with the marker for `[[ \#\# completed \#\# ]]`.
\smallskip
\end{tcolorbox}
\vspace{0.35em}

\noindent\textbf{Response}\par\smallskip

\begin{tcolorbox}[
  breakable, enhanced,
  title={Output},
  colback=OutputBg, colframe=OutputFrame,
  coltitle=white, fonttitle=\bfseries\small,
  arc=2mm, boxrule=1pt,
  before upper={\small\setlength{\parindent}{0pt}\setlength{\parskip}{2pt}},
]
\par\noindent [[ \#\# feedback \#\# ]]\\
\noindent
\par\noindent \{
\par\noindent\hangindent=1.5em\hangafter=1 \hspace*{1.5em}``\textbf{\texttt{main\_diagnosis}}\textquotesingle\textquotesingle: ``The primary failure is a systematically over-hyperpolarized potassium reversal potential (E\_K = -107 mV, hardcoded), which is far below the physiological range for cortical neurons (\textasciitilde{}-90 mV). This single constant distorts the after-hyperpolarization depth for both the standard K+ current and the M-current, biasing spike count, mean stimulation voltage, skewness, and kurtosis simultaneously. A secondary issue is that the Ih gating parameters allow unrealistically strong activation near rest: with V\_half\_r = -80 mV and k\_r = 7 mV, \textasciitilde{}19\% of Ih channels are open at -70 mV rest, and the prior permits gbar\_Ih up to 120 mS/cm$^{2}$ (physiological range is 0.05--5 mS/cm$^{2}$), causing gbar\_Ih to dominate the conductance balance and push mean resting potential and voltage variance far from target.'',
\par\noindent \hspace*{1.5em}``\textbf{\texttt{issues}}\textquotesingle\textquotesingle:
\par\noindent \hspace*{1.5em}[
\par\noindent \hspace*{3.0em}\{
\par\noindent\hangindent=4.5em\hangafter=1 \hspace*{3.0em}\hspace*{1.5em}``\textbf{\texttt{description}}\textquotesingle\textquotesingle: ``E\_K = -107.0 mV is too hyperpolarized. Physiological E\_K for cortical neurons is approximately -90 mV. The excessively negative value deepens the AHP after every spike, reduces firing rate by prolonging inter-spike intervals, pulls mean stimulation-period voltage too negative (metric 4), and skews the voltage distribution toward hyperpolarized values (metrics 6--7). Because E\_K also sets the M-current reversal, this error is amplified once gbar\_KM is inferred.'',
\par\noindent\hangindent=4.5em\hangafter=1 \hspace*{3.0em}\hspace*{1.5em}``\textbf{\texttt{location}}\textquotesingle\textquotesingle: ``Line: `E\_K = -107.0  \# mV`'',
\par\noindent\hangindent=4.5em\hangafter=1 \hspace*{3.0em}\hspace*{1.5em}``\textbf{\texttt{severity}}\textquotesingle\textquotesingle: ``critical'',
\par\noindent\hangindent=4.5em\hangafter=1 \hspace*{3.0em}\hspace*{1.5em}``\textbf{\texttt{suggestion}}\textquotesingle\textquotesingle: ``Change E\_K to -90.0 mV, which is the standard physiological value for cortical/hippocampal neurons used in most HH extensions (e.g., Pospischil et al. 2008). If the target data suggests a slightly different value, keep it as a fixed constant in the range [-80, -95] mV rather than leaving it at -107.''
\par\noindent \hspace*{3.0em}\},
\par\noindent \hspace*{3.0em}\{
\par\noindent\hangindent=4.5em\hangafter=1 \hspace*{3.0em}\hspace*{1.5em}``\textbf{\texttt{description}}\textquotesingle\textquotesingle: ``The Ih current is over-powered at rest. Two compounding problems: (1) k\_r = 7 mV with V\_half\_r = -80 mV means \textasciitilde{}19\% of Ih channels are already open at a typical resting voltage of -70 mV, providing significant tonic depolarizing current even without strong hyperpolarisation. (2) The prior upper bound for gbar\_Ih (120 mS/cm$^{2}$) is \textasciitilde{}25--100$\times$ the physiological range (0.05--5 mS/cm$^{2}$), so the inference can select values that make Ih the dominant conductance, dramatically raising resting potential SD and variance (metrics 3, 5) to unrealistic levels while also distorting skewness and kurtosis.'',
\par\noindent\hangindent=4.5em\hangafter=1 \hspace*{3.0em}\hspace*{1.5em}``\textbf{\texttt{location}}\textquotesingle\textquotesingle: ``Lines: `gbar\_Ih = params[:, 7]`, `k\_r = 7.0`, `V\_half\_r = -80.0`'',
\par\noindent\hangindent=4.5em\hangafter=1 \hspace*{3.0em}\hspace*{1.5em}``\textbf{\texttt{severity}}\textquotesingle\textquotesingle: ``major'',
\par\noindent\hangindent=4.5em\hangafter=1 \hspace*{3.0em}\hspace*{1.5em}``\textbf{\texttt{suggestion}}\textquotesingle\textquotesingle: ``Shift the half-activation to V\_half\_r = -90 mV (deeper hyperpolarisation required to open channels), which reduces Ih activation at rest from \textasciitilde{}19\% to \textasciitilde{}4\%. Also reduce the prior upper bound for gbar\_Ih from 120 to 5 mS/cm$^{2}$ to keep Ih in the physiological regime. Alternatively, increase k\_r to 10--12 mV to soften the transition and reduce resting activation. These changes together will prevent Ih from dominating resting conductance while preserving its role in correcting subthreshold distribution shape.''
\par\noindent \hspace*{3.0em}\}
\par\noindent \hspace*{1.5em}]
\par\noindent \}
\smallskip
\par\noindent [[ \#\# completed \#\# ]]\\
\noindent
\smallskip
\end{tcolorbox}
\vspace{0.35em}
